\documentclass[10pt,a4paper]{amsart}
\usepackage[margin=19mm]{geometry}
\usepackage{amsmath,amssymb,mathtools}
\usepackage[scr=rsfs,cal=euler]{mathalfa}
\usepackage{xfrac}
\usepackage[unicode,hidelinks,bookmarks=false]{hyperref}
\newcommand{\ie}{i.e.\ }
\newcommand{\cf}{cf.\ }
\newcommand{\resp}{respectively\ }
\newcommand{\Subsection}{\subsection}
\providecommand{\nobreakdash}{\nobreak\hbox{-}}
\setlength{\emergencystretch}{2em}
\multlinegap0pt
\usepackage{xcolor}
\usepackage{enumerate}
\usepackage{amssymb}
\usepackage{bm}
\newtheorem{lemma}{Lemma}
\title{Bernstein-type theorem for constant mean curvature surfaces in the isotropic 3-space}
\author{Shintaro Akamine, Wonjoo Lee, Seong-Deog Yang}
\date{Source: arXiv:2505.24109v2 (CC BY 4.0)}
\begin{document}
\maketitle

\noindent\emph{Source and licence.} Bernstein-type theorem for constant mean curvature surfaces in the isotropic 3-space. Version arXiv:2505.24109v2 (CC BY 4.0), distributed by arXiv under the Creative Commons Attribution 4.0 International licence, http://creativecommons.org/licenses/by/4.0/, as recorded in the arXiv licence metadata for this exact version and shown on the abstract page https://arxiv.org/abs/2505.24109v2. Full licence notice: \texttt{licenses/arxiv-2505.24109-CC-BY-4.0.txt}.\par\smallskip

\noindent\textit{Editorial scope.} Counted unit: the article's lemma lemma2 (source0.tex lines 255--261). The article's complete original proof of it is delivered with it (source0.tex lines 263--298, one \texttt{proof} environment of 36 lines). No mathematics is written, repaired or re-derived: the statement and the proof are the article's own lines, in the article's own notation, and no retained line is substituted or re-flowed.  Retained uncounted as the source-local context the proof needs: lines 229--231; lines 247--253.  Nothing was omitted from any retained span, so the package carries no omission record.  No retained line contains a citation, so the wrapper carries no bibliography and no bibliography entry.  No retained line contains a figure environment, an image-inclusion command or drawing code, so no image is delivered and the package declares no asset dependency.  Editorial formatting only, in addition: an AMS \texttt{amsart} preamble that restates the article's own declarations and macros used by the retained lines, namely the environments \texttt{lemma} on separate counters, together with the standard packages those lines need.  The retained environments print in this document as Lemma 1 in order of appearance, whereas the article prints the same items with the numbering of its own full document.  The complete native source of the article as distributed in the arXiv e-print for this exact version is bundled unchanged as \texttt{sources/179029/source0.tex}.
	\begin{equation}\label{eq:Wformula}
	X(z)= \Re \int \left(\overline{h_1}+h_2, 1, -i\right)\omega, \quad h_1:=H\int \omega,
	\end{equation}

\begin{lemma}\label{lemma}
The Gaussian curvature $K$ and the mean curvature $H$ of a constant mean curvature $H$ surface in $\mathbb{I}^3$ satisfy the following relation

	\begin{equation}\label{eq:KH}
	K=H^2-\left| \frac{\operatorname{d}\!h_2}{\omega} \right|^2.
	\end{equation}
\end{lemma}

\begin{lemma}\label{lemma2}
Any constant mean curvature $H$ surface with constant Gaussian curvature $K$ is a part of the surface
	\begin{equation}\label{eq:EHP}
	\ell = f(x,y)=H\left( \frac{x^2+y^2}{2}\right)+ \sqrt{H^2-K}\left( \frac{x^2-y^2}{2}\right)
	\end{equation}
up to an isometry of $\mathbb{I}^3$.
\end{lemma}

\begin{proof}
By Lemma \ref{lemma}, 
	\[
		\left|\frac{\operatorname{d}\!h_2}{\omega} \right|^2 = H^2-K.
	\]
Since the function $\operatorname{d}\!h_2/\omega $ is holomorphic, there exists a real number $\theta$ such that 
	\begin{equation}\label{eq:constancy}
	\frac{\operatorname{d}\!h_2}{\omega} =e^{i\theta}\sqrt{H^2-K}.
	\end{equation}

On the other hand, by the formula \eqref{eq:Wformula} we have the relation
	\[
		\int{\omega} = x+iy.
	\]
The relation $\hat{\omega}\neq 0$ implies that $x,y$ coordinates in $\mathbb{I}^3$ give local isothermal coordinates of the surface. Hence, we can take $z=x+iy$ as a complex coordinate so that $\omega=\operatorname{d}\!z$. 
For such a complex coordinate, $h_1$ and $h_2$ can be computed as
	\[
		h_1=Hz,\quad h_2=\sqrt{H^2-K}e^{i\theta}z
	\]
up to additive complex constants by \eqref{eq:Wformula} and \eqref{eq:constancy}.
Therefore, by using the representation formula \eqref{eq:Wformula} again, the coordinate $\ell$ of the surface can be computed as follows :
\begin{align*}
\ell &= \Re{\int{\overline{h_1}\omega}} + \Re{\int{h_2\omega}} \\ 
&=H\left(\frac{x^2+y^2}{2} \right) +\sqrt{H^2-K}\left[ \cos{\theta}\left(\frac{x^2-y^2}{2}\right)-\sin{\theta} xy\right].
\end{align*}
If we introduced new coordinates $(\tilde{x},\tilde{y})$ by the following rotation of the $xy$-plane :
	\[
		\begin{pmatrix} x \\ y \end{pmatrix} = \begin{pmatrix} \cos{\tilde{\theta}}  & -\sin{\tilde{\theta}} \\ \sin{\tilde{\theta}} & \cos{\tilde{\theta}} \end{pmatrix}
\begin{pmatrix} \tilde{x} \\ \tilde{y} \end{pmatrix}, \quad \tilde{\theta}:=-\frac{\theta}{2},
	\]
$\ell=\ell(x,y)$ can be written as
	\[
		\ell(\tilde{x},\tilde{y})=H\left(\frac{\tilde{x}^2+\tilde{y}^2}{2}\right)+\sqrt{H^2-K}\left(\frac{\tilde{x}^2-\tilde{y}^2}{2}\right)
	\]
which proves the assertion.
\end{proof}

\end{document}
